\section{Introduction}\label{sec:short-intro}
Standards may be formally incompatible but privately interoperable. A foreign producer can incur a setup cost to support another country's technical interface or certification, reducing a government's ability to exclude the producer without changing the formal membership of a regional agreement. The institutional response is not obvious: private workarounds raise member-market consumer surplus but also erode domestic producers' competitive rents, while the members' reciprocal disadvantage in an outsider market may persist.

We ask when a privately chosen, \emph{one-way} adaptation changes a regional member government's ranking of regional and international standardization. We distinguish the private adoption equilibrium from the formal institutional outcome. In a three-country Cournot environment, the pre-adaptation regional-minus-international welfare gap is $\mathscr E-\mathscr D$: a net member-market advantage less an outsider-market disadvantage. After outsider adaptation it is $-\mathscr D+R(d,v)$, where $R(d,v)$ is the residual member-market rent derived from market equilibrium. Depending on the residual cost, private adaptation can reverse the preference, strengthen the incentive without reversal, or weaken it relative to no adoption. An exact joint witness ensures that the private action generating the political effect really occurs in equilibrium.

Private converters and one-way compatibility are established firm choices \citep{farrellSaloner1992,manentiSomma2008}. \citet{gandalShy2001} show that conversion costs can sustain standardization unions, whereas strong network effects favor mutual recognition. \citet{takaradaEtAl2020} find multilateral agreements in the core for sufficiently local consumption externalities and regional agreements for large transboundary externalities; their fixed-cost extension changes regime payoffs. Here the additional-standard investment is instead chosen after the formal partition, changing the continuation value of the same union. \citet{innocentiMenicucci2021} characterize firms' compatibility choices before price competition, and \citet{dingKoShen2022} study coalition-proof partial compatibility among private platforms. Our equilibrium object is a national government's ranking of formal partitions after firms' adoption decisions. \citet{schmidtSteingress2022} connect costly firm adoption to trade responses to harmonized standards; we ask how the resulting market rents affect incentives to change a formal agreement.

A working paper by Maggi and Mr{\'{a}}zov{\'{a}} (NBER Working Paper 33318, 2024) argues that regulatory harmonization may arise without an agreement and that cooperation may instead preserve efficient diversity. Our contribution is an explicit, \emph{equilibrium-feasible} mapping from private post-agreement adaptation to \emph{national} incentives over formal standards regimes, together with failure boundaries. The coalition-stability result is a full-adaptation application, not a theorem about incomplete-adaptation stability.

\paragraph{Relation to public policy and private standards choice} \citet{fischerSerra2000} and \citet{grossmanMcCalmanStaiger2021} study the costs of serving markets governed by different standards, while \citet{klimenko2009} analyzes public policies and international agreements concerning technical compatibility. \citet{kawabataTakarada2021} examine whether regional standards cooperation impedes wider harmonization. On the firm side, \citet{buccellaEtAl2023} study endogenous compatibility with network effects, and \citet{parentiVannoorenberghe2025} link costly local product tailoring to regulatory blocs. Regulatory product standards can address consumption externalities, whereas compatibility here changes demand and producer cost. A different government objective alone is not our contribution: the model derives the national welfare effect of an outsider's private adoption after a formal partition is selected, including the firm continuation that makes that adoption feasible.

Section~\ref{sec:short-model} sets out the primitives. Sections~\ref{sec:short-adoption} and \ref{sec:short-incentives} solve private adoption and government ranking. Section~\ref{sec:short-stability} examines institutional stability and robustness. Section~\ref{sec:short-discussion} discusses interpretation and limits.
